\chapter{Architecture}

This chapter defines the v2 system architecture: the five primary components, their interfaces and invariants,
the complete data flow from runbook invocation to trace archive, the execution lifecycle, the concurrency model,
and how the \texttt{gert serve} JSON-RPC layer integrates as a first-class adapter.

The single most important architectural rule is: \textbf{dependencies flow inward toward the Core Domain.}
The Runtime Core knows nothing about TUI, VS Code, or JSON-RPC. Adapters know nothing about execution semantics.
Every component boundary is enforced by a Go interface; no component imports a package that is ``above'' it in
the dependency graph.

% ─────────────────────────────────────────────────────────────────────────────
\section{Primary Components}
% ─────────────────────────────────────────────────────────────────────────────

\subsection{Parser}

\paragraph{Responsibility.}
The Parser owns the transformation of a raw YAML document into a validated, semantically-resolved
\texttt{ParsedRunbook} value. It does \emph{not} own execution, planning, tool discovery, or input resolution.
It validates schema structure and surface-level semantic rules (duplicate step IDs, unknown step types,
malformed Go template expressions) but intentionally defers deep semantic validation (e.g., whether a
\texttt{from:} binding can actually be resolved) to later pipeline stages.

\paragraph{Interface.}

\begin{verbatim}
// Parser transforms raw YAML bytes into a validated ParsedRunbook.
type Parser interface {
    // Parse reads YAML from r, validates against the v2 JSON Schema,
    // applies structural semantic checks, and returns a ParsedRunbook.
    // Errors carry structured diagnostics (file, line, column, message).
    Parse(ctx context.Context, r io.Reader, opts ParseOptions) (*ParsedRunbook, error)
}

type ParseOptions struct {
    // SchemaVersion overrides the apiVersion field for compatibility shims.
    SchemaVersion string
    // StrictMode causes deprecation warnings to be surfaced as errors.
    StrictMode bool
}

// ParsedRunbook is an immutable, fully-validated representation of a runbook.
// All fields have been structurally validated. Templates are not yet evaluated.
type ParsedRunbook struct {
    Meta       RunbookMeta
    Governance GovernancePolicy  // nil if no governance block
    Steps      []ParsedStep
    Imports    map[string]string // alias -> path (not yet resolved)
    Tools      []string          // tool names (not yet discovered)
    Schema     SchemaRef         // version and $schema URL
}
\end{verbatim}

\paragraph{Dependencies.}
The Parser depends only on the JSON Schema artifact (a static embedded file) and the expression
evaluator (for template syntax validation). It has no runtime dependencies.

\paragraph{Key Invariants.}
\begin{itemize}
  \item A \texttt{ParsedRunbook} returned without error is structurally valid against the v2 JSON Schema.
  \item All step IDs are unique within a \texttt{ParsedRunbook}.
  \item No I/O occurs during parsing (file imports are recorded as paths, not resolved).
  \item Parsing is pure and deterministic: the same input always produces the same output.
\end{itemize}

% ──────────────────────────────────────────────────────────────────────────────
\subsection{Planner}

\paragraph{Responsibility.}
The Planner transforms a \texttt{ParsedRunbook} into an \texttt{ExecutionPlan}: a resolved, ordered,
dependency-annotated representation of every step the runtime will execute. The Planner resolves imports
(recursively parsing and inlining invoked runbooks), discovers and validates tool definitions, resolves
input provider bindings to their declared sources, and produces the static step ordering. It does
\emph{not} evaluate dynamic conditions (branch predicates, iterate count expressions) — those are
runtime concerns.

\paragraph{What is an ExecutionPlan?}
An \texttt{ExecutionPlan} is a \emph{flat, ordered list} of resolved steps with pre-computed metadata,
not a DAG. Branch logic and iteration are left as runtime decisions because they depend on captured
variables evaluated during execution. The plan encodes \emph{which steps exist} and \emph{what
preconditions they carry}, not which steps will actually run. This design keeps the planner pure and
the runtime authoritative over control flow.

\begin{verbatim}
type Planner interface {
    // Plan resolves imports, discovers tools, validates bindings,
    // and returns a fully-resolved ExecutionPlan ready for the runtime.
    Plan(ctx context.Context, rb *ParsedRunbook, opts PlanOptions) (*ExecutionPlan, error)
}

type PlanOptions struct {
    BaseDir     string            // working directory for import resolution
    ToolDirs    []string          // directories to search for .tool.yaml files
    ProviderDir string            // directory for .provider.yaml files
    Vars        map[string]string // CLI-supplied variable overrides
}

// ExecutionPlan is an immutable, fully-resolved plan for a single run.
type ExecutionPlan struct {
    RunID       string
    RunbookPath string
    Steps       []ResolvedStep       // ordered; includes inlined invoke steps
    Tools       map[string]*ToolDef  // name -> resolved tool definition
    Providers   map[string]*ProviderDef // binding name -> provider
    Governance  *GovernancePolicy
    Metadata    PlanMetadata
}

// ResolvedStep is a single step in the plan, fully resolved.
type ResolvedStep struct {
    ID       string
    Kind     StepKind      // cli, manual, tool, invoke, branch, iterate
    Spec     StepSpec      // kind-specific parameters (resolved, not templated)
    Contract *contract.Contract  // effects/reads/writes/idempotent/deterministic
    Depth    int           // chain depth (0 = root runbook)
    Origin   string        // source runbook path (for inlined invoke steps)
}
\end{verbatim}

\paragraph{Dependencies.}
The Planner depends on the Parser (for import resolution), the Tool Registry (for
\texttt{.tool.yaml} discovery), and the Provider Registry (for \texttt{.provider.yaml}
discovery). It must not depend on the Runtime Core, Extension Host, or any adapter.

\paragraph{Key Invariants.}
\begin{itemize}
  \item An \texttt{ExecutionPlan} returned without error contains only resolvable steps.
  \item All tool names referenced by steps appear in \texttt{ExecutionPlan.Tools}.
  \item Import graphs are acyclic; the Planner detects and rejects cycles.
  \item All \texttt{from:} bindings that declare a provider appear in \texttt{ExecutionPlan.Providers}.
  \item The Planner does not execute or side-effect; it is safe to plan without running.
\end{itemize}

% ──────────────────────────────────────────────────────────────────────────────
\subsection{Runtime Core}

\paragraph{Responsibility.}
The Runtime Core owns the execution loop: advancing through an \texttt{ExecutionPlan} step by step,
evaluating templates and branch predicates, dispatching tool invocations, collecting evidence, writing
the append-only trace, enforcing governance pre-flight checks, emitting structured events, and
managing the mutable run state. It does \emph{not} own user presentation, serialization formats, or
network transport.

\paragraph{Interface.}

\begin{verbatim}
type Runtime interface {
    // Start initializes a run from a plan and returns a RunHandle.
    // The run does not advance until Next is called.
    Start(ctx context.Context, plan *ExecutionPlan, opts RunOptions) (RunHandle, error)

    // Resume restores a previously checkpointed run from the run store.
    Resume(ctx context.Context, runID string, opts RunOptions) (RunHandle, error)
}

type RunHandle interface {
    // Next advances the run by one step. Blocks until the step completes
    // (or is paused waiting for evidence/approval).
    // Returns io.EOF when the run has no more steps.
    Next(ctx context.Context) (*StepResult, error)

    // Approve records approval for the current step's approval gate.
    Approve(ctx context.Context, decision ApprovalDecision) error

    // SubmitEvidence records evidence for the current manual step.
    SubmitEvidence(ctx context.Context, stepID string, ev map[string]*EvidenceValue) error

    // Cancel requests a graceful stop of the run.
    Cancel(ctx context.Context, reason string) error

    // State returns the current mutable run state (read-only snapshot).
    State() RunState

    // Events returns a channel of structured events emitted during execution.
    Events() <-chan Event
}

type RunOptions struct {
    Mode        RunMode           // real, dry-run, replay
    Actor       string            // --as identity
    Vars        map[string]string // variable overrides
    ScenarioDir string            // replay scenario directory
    Store       RunStore          // durable store for checkpointing
    OnEvent     func(Event)       // synchronous event hook (nil = use Events channel)
}
\end{verbatim}

\paragraph{Dependencies.}
The Runtime Core depends on the Governance Engine, Tool Runtime (for tool invocation),
Expression Evaluator (for templates and branch predicates), Evidence store, and Trace Writer.
It does not depend on the Extension Host, API/Adapter Layer, or any adapter.

\paragraph{Key Invariants.}
\begin{itemize}
  \item The trace is append-only: no event is ever deleted or modified after being written.
  \item Governance pre-flight is enforced before every \texttt{cli} and \texttt{tool} step dispatch.
  \item Redaction is applied to all captured output before it is written to the trace.
  \item The run state is checkpointed after every step completion.
  \item Events are emitted in a total order that matches step execution order.
  \item The runtime is re-entrant per run: multiple callers may call \texttt{Next} sequentially
        (but not concurrently) on the same \texttt{RunHandle}.
\end{itemize}

% ──────────────────────────────────────────────────────────────────────────────
\subsection{Extension Host}

\paragraph{Responsibility.}
The Extension Host manages the lifecycle of out-of-process extensions: discovery, launch, capability
negotiation, RPC dispatch, sandboxing, and graceful shutdown. It exposes extension-contributed tools
and providers to the Tool Runtime and Provider Registry. It does \emph{not} execute runbook steps
directly and does not own the extension communication protocol (that belongs to the extension RPC
layer, defined in §04).

\paragraph{Interface.}

\begin{verbatim}
type ExtensionHost interface {
    // Load discovers and launches all extensions declared in the project manifest.
    Load(ctx context.Context, manifest *ProjectManifest) error

    // Shutdown gracefully stops all running extensions.
    Shutdown(ctx context.Context) error

    // ContributedTools returns all tools contributed by loaded extensions.
    ContributedTools() []*ToolDef

    // ContributedProviders returns all providers contributed by loaded extensions.
    ContributedProviders() []*ProviderDef

    // ContributedPolicyRules returns governance rules contributed by extensions.
    // These compose with host policy; host policy takes precedence on conflict.
    ContributedPolicyRules() []PolicyRule
}
\end{verbatim}

\paragraph{Dependencies.}
The Extension Host depends on the extension manifest schema and the extension RPC transport
(stdio JSON-RPC, as specified in §04). It does not depend on the Runtime Core or any adapter.

\paragraph{Key Invariants.}
\begin{itemize}
  \item No extension operates without an explicit capability grant from the host.
  \item Extension crashes do not crash the host process.
  \item The Extension Host enforces timeout and cancellation on all RPC calls to extensions.
  \item Extension-contributed policy rules are additive; they can only \emph{restrict}, not
        \emph{relax}, host-defined policy.
\end{itemize}

% ──────────────────────────────────────────────────────────────────────────────
\subsection{API / Adapter Layer}

\paragraph{Responsibility.}
The API/Adapter Layer is the boundary between the Runtime Core and all external consumers
(TUI, VS Code extension, web clients, CI pipelines). It translates the runtime's event stream
and \texttt{RunHandle} contract into transport-specific protocols. Each adapter is a thin,
stateless renderer: it subscribes to the event stream, presents state to the user, and forwards
user input (approvals, evidence, cancellation) back to the \texttt{RunHandle}. Adapters own
\emph{no} execution logic.

\paragraph{Interface.}

\begin{verbatim}
// Adapter is the interface every consumer of the Runtime Core must implement.
type Adapter interface {
    // Attach binds the adapter to a RunHandle before the first Next call.
    Attach(handle RunHandle) error

    // Run blocks until the run completes or is cancelled. It is responsible
    // for draining Events(), rendering output, and forwarding user input.
    Run(ctx context.Context) error
}

// The gert serve JSON-RPC layer implements Adapter over stdio.
// The TUI (Bubble Tea) implements Adapter over a terminal.
// Future: a web adapter implements Adapter over WebSocket.
\end{verbatim}

\paragraph{Dependencies.}
Adapters depend on the \texttt{RunHandle} and \texttt{Event} types from the Runtime Core.
They must not import any package from the Parser, Planner, or Extension Host.

\paragraph{Key Invariants.}
\begin{itemize}
  \item An adapter may not call \texttt{Next} more than once concurrently.
  \item An adapter may not modify run state directly; all state changes go through \texttt{RunHandle}.
  \item Adapters must handle the case where the event channel is closed (run ended) without blocking.
  \item An adapter failure must not corrupt the run trace or run state.
\end{itemize}

\paragraph{Dependency Direction Summary.}

\begin{verbatim}
  Adapter Layer
       |
       v
  Runtime Core  <---- Extension Host
       |                    |
       v                    v
    Planner           Tool Runtime
       |
       v
    Parser
       |
       v
  Schema / JSON Schema artifact
\end{verbatim}

No upward arrow exists. The Parser knows nothing about the Runtime. The Runtime knows nothing
about adapters. The Extension Host is a peer of the Runtime Core, contributing tools and policy
rules but not owning execution.

% ─────────────────────────────────────────────────────────────────────────────
\section{Data Flow}
% ─────────────────────────────────────────────────────────────────────────────

This section traces the complete path from runbook invocation to trace archive, identifying
exactly where each responsibility is exercised.

\subsection{Phase 1: Parse}

The user invokes \texttt{gert exec runbook.yaml --as eng@corp.com}. The CLI layer reads the
YAML file and calls \texttt{Parser.Parse()}. The Parser validates the document against the
embedded v2 JSON Schema, checks structural semantic rules, and returns a \texttt{ParsedRunbook}.
If validation fails, structured diagnostics are returned and execution does not proceed.

\subsection{Phase 2: Plan}

The \texttt{ParsedRunbook} is passed to \texttt{Planner.Plan()}. The Planner:
\begin{enumerate}
  \item Resolves all \texttt{imports:} recursively (parse → plan each imported runbook, inline
        its steps with incremented \texttt{Depth}).
  \item Discovers \texttt{.tool.yaml} files for each named tool.
  \item Resolves \texttt{from:} input bindings to their declared providers.
  \item Validates that all step references (branch targets, iterate bodies) are valid step IDs.
  \item Constructs and returns an \texttt{ExecutionPlan}.
\end{enumerate}

Governance policy from the runbook's \texttt{meta.governance} block is extracted and embedded
in the \texttt{ExecutionPlan} for runtime enforcement.

\subsection{Phase 3: Run Initialization}

\texttt{Runtime.Start(plan, opts)} is called. The runtime:
\begin{enumerate}
  \item Generates a \texttt{RunID} (format: \texttt{YYYYMMDDTHHmmss-xxxxxxxx}).
  \item Creates the run directory: \texttt{.runbook/runs/<runID>/}.
  \item Opens the append-only trace file: \texttt{.runbook/runs/<runID>/trace.jsonl}.
  \item Writes the \texttt{run/started} trace event (runbook path, actor, mode, plan hash).
  \item Initializes the \texttt{GovernanceEngine} from the plan's governance policy.
  \item Compiles redaction rules from the governance policy.
  \item Returns a \texttt{RunHandle}; the run is now paused at step 0.
\end{enumerate}

\subsection{Phase 4: Step Execution Loop}

The adapter calls \texttt{RunHandle.Next(ctx)} in a loop until \texttt{io.EOF}.
For each step, the runtime executes the following sequence:

\subsubsection*{Step Pre-flight (all step types)}
\begin{enumerate}
  \item Write \texttt{step/started} trace event.
  \item Emit \texttt{StepStarted} runtime event to the event channel.
  \item Evaluate template expressions in the step spec against current run variables.
\end{enumerate}

\subsubsection*{Governance Pre-flight (cli and tool steps)}
\begin{enumerate}
  \item \textbf{Command check:} \texttt{GovernanceEngine.CheckCommand(argv[0])} evaluates
        the denylist (deny takes precedence) and then the allowlist. A violation writes a
        \texttt{governance/commandBlocked} trace event and returns an error; the step fails.
  \item \textbf{Env var check:} \texttt{GovernanceEngine.FilterEnvVars()} removes variables
        matching \texttt{deny\_env\_vars} patterns. Blocked variable names are written to the
        \texttt{governance/envVarBlocked} trace event.
  \item \textbf{Contract check:} The step's \texttt{Contract.Risk()} is evaluated. Steps
        with \texttt{RiskCritical} that do not have an explicit approval gate declaration
        in the plan emit a \texttt{governance/contractWarning} event (not a hard failure in
        v2.0; upgraded to error in strict mode).
  \item \textbf{Approval gate:} If the step declares \texttt{approvals.min > 0}, the runtime
        pauses and emits an \texttt{ApprovalRequired} event. Execution blocks at this step
        until \texttt{RunHandle.Approve()} is called with sufficient approvals.
        See §\ref{sec:approval-gate} (Governance chapter) for the full approval protocol.
\end{enumerate}

\subsubsection*{Dispatch (step-type-specific)}
\begin{itemize}
  \item \textbf{cli:} Fork a subprocess. Capture stdout/stderr. Apply redaction rules to
        captured output before storing in run variables or writing to trace.
  \item \textbf{tool:} Call \texttt{ToolRuntime.Invoke(toolName, input)}. The tool runtime
        dispatches to the appropriate transport (stdio, JSON-RPC, MCP). Redaction applied
        to returned output.
  \item \textbf{manual:} Emit \texttt{EvidenceRequired} event. Block until
        \texttt{SubmitEvidence()} is called. SHA256-hash each evidence attachment.
  \item \textbf{invoke:} Inlined at plan time; at runtime the invoke body steps are already
        present in the plan at \texttt{Depth+1}. No special dispatch required.
  \item \textbf{branch:} Evaluate branch predicates (expression evaluator) against current
        variables. Select the matching branch arm. Write \texttt{event/branchResolved}
        trace event. Advance step cursor to the first step of the resolved arm.
  \item \textbf{iterate:} Evaluate the iteration expression. Write
        \texttt{event/iteratePassStart}. Execute body steps. Write
        \texttt{event/iteratePassEnd}. Repeat until condition is false.
\end{itemize}

\subsubsection*{Step Post-flight (all step types)}
\begin{enumerate}
  \item Apply outcome evaluation: test \texttt{when:} expressions; record \texttt{OutcomeRecord}.
  \item Write captured variables to run state.
  \item Checkpoint run state to \texttt{.runbook/runs/<runID>/snapshots/<stepID>.json}.
  \item Write \texttt{step/completed} trace event (outcome, captures, elapsed time).
  \item Emit \texttt{StepCompleted} runtime event.
\end{enumerate}

\subsection{Phase 5: Run Completion}

When \texttt{Next} returns \texttt{io.EOF}:
\begin{enumerate}
  \item Write \texttt{run/completed} trace event (final outcome, step counts, elapsed time).
  \item Flush and sync the trace file.
  \item Compute SHA256 of the trace file; write \texttt{run/evidenceHash} record.
  \item Optionally archive the run directory (compress to \texttt{.runbook/archive/}).
  \item Close the event channel.
\end{enumerate}

% ─────────────────────────────────────────────────────────────────────────────
\section{Lifecycle Model}
% ─────────────────────────────────────────────────────────────────────────────

\subsection{Run Start}

A run starts in one of two ways:

\begin{enumerate}
  \item \textbf{Direct invocation:} The CLI or test harness calls \texttt{Runtime.Start()}
        directly. This is the path for \texttt{gert exec}, \texttt{gert tui}, \texttt{gert debug},
        and \texttt{gert test}.
  \item \textbf{RPC invocation:} A client sends \texttt{exec/start} to the JSON-RPC server
        (\texttt{gert serve}). The server calls \texttt{Runtime.Start()} on behalf of the
        client and returns the \texttt{RunID} in the response. Subsequent step advancement
        is driven by \texttt{exec/next} RPC calls.
\end{enumerate}

\subsection{Step Advancement}

Step advancement is \emph{always client-driven}. The runtime does not auto-advance.
After \texttt{Start}, the client calls \texttt{Next} to execute each step in sequence.
This applies to all adapters, including headless CI mode. The only exception is the
\texttt{--auto} flag, which causes the CLI adapter to call \texttt{Next} in a tight loop,
producing the appearance of automatic advancement while retaining the single-step contract.

\subsection{Pause and Resume}

A run is \emph{paused} when \texttt{Next} returns a non-terminal result that requires external
input: evidence submission (\texttt{manual} step) or approval (\texttt{approval} gate).
The run remains paused until the required input is provided via \texttt{SubmitEvidence} or
\texttt{Approve}. The adapter is responsible for presenting the pause to the user.

The checkpoint written after each completed step is sufficient to resume a run after a
process restart. The checkpoint records: completed step IDs, current variable bindings,
current evidence values, and the trace sequence number. Resumption is initiated by calling
\texttt{Runtime.Resume(runID)} from any adapter.

\subsection{Cancellation}

\texttt{RunHandle.Cancel(ctx, reason)} initiates graceful cancellation:
\begin{enumerate}
  \item The currently-executing subprocess (if any) receives \texttt{SIGTERM} followed by
        \texttt{SIGKILL} after a 5-second grace period.
  \item The runtime writes a \texttt{run/cancelled} trace event with the reason and actor.
  \item Subsequent \texttt{Next} calls return \texttt{ErrRunCancelled}.
  \item The run state is checkpointed as \texttt{cancelled} (resumable).
\end{enumerate}

The cancellation contract does not invoke compensation handlers in v2.0. Saga/compensation
(stepping backward through a compensation chain on failure) is deferred to v2.1 as specified
in §09 (Open Questions).

\subsection{Crash Recovery}

If the gert process crashes mid-run, the trace and checkpoint files remain on disk in a
consistent state due to the append-only, sync-after-write trace discipline. On restart,
the user (or the VS Code extension) can call \texttt{Runtime.Resume(runID)} to continue
from the last completed step. The runtime reads the checkpoint to reconstruct run state.
Steps already recorded as completed in the trace are not re-executed.

% ─────────────────────────────────────────────────────────────────────────────
\section{Concurrency Model}
% ─────────────────────────────────────────────────────────────────────────────

\subsection{Per-Process Run Isolation}

The v2 runtime is \textbf{single-run-per-process}. Each \texttt{gert exec} invocation is a
separate OS process. \texttt{gert serve} is the exception: it is a long-running process that
may host multiple sequential runs (one at a time) for a single VS Code extension session.
Concurrent runs within a single \texttt{gert serve} instance are explicitly out of scope
for v2.0 (see §09).

\paragraph{Rationale.} Governance enforcement, file-level trace writes, and subprocess management
are dramatically simpler with one run per process. Process isolation also provides a natural
crash boundary: a failed run cannot corrupt the state of a concurrent run.

\subsection{Goroutine Model}

Within a single run, the concurrency model is:

\begin{itemize}
  \item \textbf{One goroutine per run:} the step execution loop runs on a single goroutine.
        Steps execute sequentially; there is no goroutine-per-step parallelism in v2.0.
  \item \textbf{Event fan-out goroutine:} a second goroutine reads from an internal event
        buffer and writes to the \texttt{Events()} channel. This decouples event emission
        (which must be synchronous with step execution) from event consumption (which may
        be slow, e.g., a WebSocket write).
  \item \textbf{Subprocess goroutine:} when a \texttt{cli} step spawns a subprocess, a
        goroutine is used to drain stdout/stderr concurrently, preventing deadlock on
        full pipe buffers.
  \item \textbf{Tool RPC goroutine:} tool invocations over JSON-RPC or MCP use a goroutine
        per call, managed by the Tool Runtime's connection pool.
\end{itemize}

\subsection{Thread Safety at Component Boundaries}

\begin{itemize}
  \item \texttt{RunHandle} is \textbf{not safe for concurrent use}. The adapter must
        serialize all calls. \texttt{gert serve} enforces this with a per-session mutex.
  \item The \texttt{Events()} channel is safe for a single consumer goroutine.
  \item The \texttt{GovernanceEngine} is \textbf{read-only after construction} and therefore
        safe for concurrent read access (relevant if future versions support parallel steps).
  \item The Trace Writer is protected by an internal mutex; write calls are safe to issue
        from any goroutine within the run.
\end{itemize}

% ─────────────────────────────────────────────────────────────────────────────
\section{\texttt{gert serve} Integration}
% ─────────────────────────────────────────────────────────────────────────────

\texttt{gert serve} is the JSON-RPC 2.0 server that powers the VS Code extension, the MCP
tools layer, and any future web client. In the v2 architecture it is classified as an
\textbf{adapter in the API/Adapter Layer}.

\subsection{Architectural Position}

\texttt{gert serve} is a process that:
\begin{enumerate}
  \item Starts a long-lived stdio JSON-RPC 2.0 listener (newline-delimited JSON over stdin/stdout).
  \item On receiving \texttt{exec/start}, calls \texttt{Runtime.Start()} and stores the
        resulting \texttt{RunHandle} in a session map (keyed by \texttt{RunID}).
  \item On receiving \texttt{exec/next}, calls \texttt{RunHandle.Next()} and returns the
        \texttt{StepResult} as the RPC response.
  \item Drains the \texttt{Events()} channel in a background goroutine and forwards each
        event as a JSON-RPC notification to the client.
  \item On receiving \texttt{exec/approve}, \texttt{exec/submitEvidence}, or
        \texttt{exec/cancel}, dispatches to the corresponding \texttt{RunHandle} method.
\end{enumerate}

\subsection{Event Forwarding}

Every event emitted by the Runtime Core is forwarded to the client as a JSON-RPC
\emph{notification} (no \texttt{id} field) with method \texttt{event/\textit{eventType}}.
The payload is the event's JSON representation. The VS Code extension subscribes to
these notifications to update its UI without polling.

Events forwarded over WebSocket (for a future web adapter) follow the same schema;
only the transport changes.

\subsection{Crash Recovery in \texttt{gert serve}}

If the VS Code extension reconnects after a crash (e.g., the gert process was killed),
\texttt{gert serve} exposes \texttt{exec/list} and \texttt{exec/resume} methods. The
extension discovers the interrupted run, calls \texttt{exec/resume}, and receives a new
\texttt{RunHandle} over the same session. The trace and checkpoint on disk ensure no
steps are lost.

\subsection{RPC Method Summary}

\begin{verbatim}
exec/start          → Start a new run (returns RunID, first StepResult)
exec/next           → Advance run by one step (returns StepResult)
exec/approve        → Submit approval decision for current gate
exec/submitEvidence → Submit evidence for current manual step
exec/cancel         → Cancel the active run
exec/list           → List all runs for this workspace (active + completed)
exec/resume         → Resume a previously interrupted run
runbook/validate    → Validate a runbook YAML (parse + plan, no exec)
runbook/diagram     → Generate a Mermaid or ASCII diagram
event/<type>        → Notification from server to client (not a request)
\end{verbatim}
